\documentclass[11pt]{article}
\usepackage{amsmath,amssymb}
\usepackage{physics}
\usepackage[margin=1in]{geometry}
\usepackage{enumitem}

\title{ATPL 2025 --- Exercises: Efficient Quantum Algorithms and the QFT}
\author{Lecture 1: Why do we need hybrid programming models?}
\date{}

\begin{document}
\maketitle

Background: Nielsen \& Chuang, Section 5.1 (The quantum Fourier
transform), and the "Examples of Quantum Advantage" / Shor's algorithm
discussion from Lecture 1.

\begin{enumerate}[label=\textbf{Exercise \arabic*.}, leftmargin=2.2cm]

\item \textbf{QFT matrix for small $n$.}
The $N$-point QFT acts on computational basis states as
$\mathrm{QFT}_N \ket{j} = \frac1{\sqrt N}\sum_{k=0}^{N-1}
\omega^{jk}\ket{k}$, with $\omega = e^{2\pi i/N}$.
\begin{enumerate}[label=(\alph*)]
\item Write down the $2\times2$ matrix for $\mathrm{QFT}_2$ (i.e.\ $N=2$,
one qubit) explicitly, and show it equals the Hadamard gate $H$.
\item Write down the $4\times4$ matrix for $\mathrm{QFT}_4$ (two qubits,
$N=4$, $\omega = i$) explicitly, with each entry as a power of $i$
divided by 2. Verify unitarity by checking that two different columns
are orthonormal.
\end{enumerate}

\item \textbf{QFT circuit and gate count.}
The standard QFT circuit on $n$ qubits uses one Hadamard per qubit and a
cascade of controlled phase gates $R_k = \mathrm{diag}(1,
e^{2\pi i/2^k})$ between qubit pairs, followed by a reversal of qubit
order (via SWAPs).
\begin{enumerate}[label=(\alph*)]
\item Describe (in words, or as a gate list) the $\mathrm{QFT}_3$
circuit (3 qubits) using $H$ and controlled-$R_k$ gates.
\item Count the total number of 2-qubit gates as a function of $n$, and
compare this to the $\mathcal O(4^n)$-type cost of generic state
preparation from the previous exercise sheet. Why is the QFT "cheap"
despite acting on an exponentially large state space?
\end{enumerate}

\item \textbf{QFT of a basis state.}
Write down a closed-form expression for $\mathrm{QFT}_N\ket{x}$ for
general $x \in \{0,\dots,N-1\}$ (this is just the defining sum with
$j=x$), and specialize it to $n=2$ qubits, $x=2$ (i.e.\ compute
$\mathrm{QFT}_4\ket{10}$ explicitly as a 4-component vector). Interpret
the result as a "pure frequency $x$" state.

\item \textbf{Toy period-finding.}
Consider the 3-qubit ($N=8$) equal superposition over a periodic set of
positions with period $r=2$:
\[
\ket\psi = \frac12\left(\ket0+\ket2+\ket4+\ket6\right).
\]
\begin{enumerate}[label=(\alph*)]
\item Using linearity and Exercise 3, compute $\mathrm{QFT}_8\ket\psi$
as a sum of four QFT basis-state outputs, and show that the result is
supported \emph{only} on multiples of $N/r = 4$, i.e.\ only on $\ket0$
and $\ket4$ have nonzero amplitude.
\item Explain, in one or two sentences, how measuring the QFT output
would let you recover the period $r=2$, and how this toy example
captures the essential idea (with a modular-exponentiation oracle in
place of the hand-built $\ket\psi$ above) behind the quantum part of
Shor's algorithm.
\end{enumerate}

\item \textbf{The QFT does not give you a free lunch.}
Even though the QFT circuit itself is cheap ($\mathrm{poly}(n)$ gates),
use the \emph{Output Problem} from the lecture to explain why you
cannot simply read off all $2^n$ Fourier amplitudes of an arbitrary
input state from a single run. What can you actually extract
efficiently from one (or a few) runs, and why is that enough for
Shor's algorithm to work?

\end{enumerate}

\end{document}
